\chapter{Compactness at a fixed phase-field parameter}\label{sec:phase-fixed-parameter}
Here $(M^n,g)$ is a smooth nonempty closed Riemannian manifold, $n\geq1$, and $\varepsilon>0$ is fixed. We develop the compactness calculation suggested in \cite{ChAC}, Exercise 4.4, p.~14. The residual retains the correct $\varepsilon^2$ factor. The needed Sobolev compactness is proved in Section \ref{sec:sobolev-closed-compactness}; elliptic regularity remains an analytic input identified in the proofs. The resulting compactness statement alone is not a mountain-pass existence theorem.

\begin{proposition}[An energy identity for every closed critical point]\label{prop:phase-critical-energy-identity}
A smooth solution of $\varepsilon^2\Delta_g u=u^3-u$ satisfies
\[
 \varepsilon^2\int_M|\nabla u|^2=\int_M(u^2-u^4),\qquad
 E_\varepsilon(u;M)=\frac{\operatorname{Vol}(M)-\int_Mu^4}{4\varepsilon}.
\]
In particular $0\leq E_\varepsilon(u;M)\leq\operatorname{Vol}(M)/(4\varepsilon)$ and
$\int_M|\nabla u|^2\leq\operatorname{Vol}(M)/(4\varepsilon^2)$.
Equality in the energy upper bound occurs precisely for $u=0$.
\end{proposition}
\begin{proof}
Multiply the equation by $u$ and integrate by parts on the closed manifold. This gives the first identity. Substitute its gradient term into
$E_\varepsilon=\varepsilon\int|\nabla u|^2/2+(4\varepsilon)^{-1}\int(1-2u^2+u^4)$; the $u^2$ terms cancel and give the second. Nonnegativity follows from the original energy definition. Proposition \ref{prop:phase-maximum-bound} gives $|u|\leq1$, and $u^2-u^4\leq1/4$, proving the gradient bound. The energy upper bound follows from $u^4\geq0$, and equality forces $u=0$ by continuity. The identities concern critical points, not arbitrary functions with the same pointwise bound.
\end{proof}

\begin{theorem}[Strong compactness for a bounded residual sequence]\label{thm:phase-fixed-parameter-compactness}
Let $u_j\in H^1(M)$ satisfy $|u_j|\leq1$ almost everywhere and $E_\varepsilon(u_j;M)\leq C$. Suppose that
\[
 r_j=-\varepsilon^2\Delta_g u_j+W'(u_j)\longrightarrow0
       \quad\text{in }H^{-1}(M),
\]
where $H^{-1}$ is the dual of $H^1(M)$ with its full norm. Then a subsequence converges strongly in $H^1$ and in every finite $L^p$ to a smooth critical point $u$, with $|u|\leq1$ and $E_\varepsilon(u_j;M)\to E_\varepsilon(u;M)$. This develops the bounded-range compactness mechanism in \cite{ChAC}, Exercise 4.4, with the residual normalized to the displayed $\varepsilon$-equation. For this same residual, convergence in $L^2$ implies the $H^{-1}$ convergence assumed here.
\end{theorem}
\begin{proof}
The pointwise bound controls the $L^2$ norms, and the energy controls the gradients because $\varepsilon$ is fixed. Weak $H^1$ compactness and Theorem \ref{thm:sobolev-closed-l2-compactness} supply a subsequence with $u_j\rightharpoonup u$ in $H^1$ and $u_j\to u$ strongly in $L^2$ and almost everywhere. The limit lies in $[-1,1]$. Since $W'$ is Lipschitz on this interval, $W'(u_j)\to W'(u)$ in $L^2$. Passing to the equation against smooth tests shows that $u$ is a weak solution, and density gives the equation against every $H^1$ test.

Subtract the limit equation and test with $u_j-u$:
\[
 \varepsilon^2\int|\nabla(u_j-u)|^2
 =\langle r_j,u_j-u\rangle
   -\int(W'(u_j)-W'(u))(u_j-u).
\]
The dual pairing tends to zero because the tests stay bounded in $H^1$, and the second term tends to zero by strong $L^2$ convergence. Thus convergence is strong in $H^1$. Boundedness by one upgrades it to every finite $L^p$: for $p\geq2$ use $|u_j-u|^p\leq2^{p-2}|u_j-u|^2$, and for smaller $p$ use finite volume. The potential $W$ is Lipschitz on $[-1,1]$, so its integrals converge; strong gradient convergence then gives energy convergence.

Finally a bounded weak solution of this smooth semilinear equation is smooth by interior elliptic $W^{2,p}$ estimates followed by Sobolev embedding and bootstrapping in finitely many coordinate charts. These elliptic regularity inputs are external here. There is no restriction $n\leq4$: the assumed uniform range bound supplies the integrability of the nonlinearity even when $H^1$ alone would not embed into $L^4$.
\end{proof}

\begin{corollary}[The entire critical set is compact for fixed epsilon]\label{cor:phase-fixed-critical-set}
For a fixed $\varepsilon>0$ and fixed smooth closed $(M,g)$, every sequence of smooth Allen--Cahn critical points has a subsequence converging in $C^k(M)$ for every finite $k$ to a smooth critical point. The constants in these compactness estimates may depend on $\varepsilon$.
\end{corollary}
\begin{proof}
The maximum bound and Proposition \ref{prop:phase-critical-energy-identity} supply all assumptions of Theorem \ref{thm:phase-fixed-parameter-compactness}, with residual zero. For the stronger convergence, the equation has uniformly bounded right-hand side after division by the fixed $\varepsilon^2$. Global estimates obtained from the finitely many local elliptic charts give uniform $W^{2,p}$ bounds for every finite $p$. Choose $p>n$ to control $C^{1,\alpha}$, and bootstrap the smooth polynomial nonlinearity to uniform bounds of all derivative orders. These are the same external elliptic estimates used in the theorem. Arzela--Ascoli and a diagonal subsequence give every $C^k$ convergence. The limit equation follows either from this convergence or from the preceding weak argument. These estimates do not assert compactness as $\varepsilon\downarrow0$.
\end{proof}

\begin{proposition}[Index cannot increase in a bounded-range strong limit]\label{prop:phase-index-lower-semicontinuity}
Fix $\varepsilon>0$. If $|u_j|,|u|\leq1$ almost everywhere and $u_j\to u$ in $L^2(M)$, define the index of each quadratic form
\[
 Q_{\varepsilon,u}(\phi)=\int_M
       \bigl(\varepsilon|\nabla\phi|^2+\varepsilon^{-1}(3u^2-1)\phi^2\bigr)
\]
as the supremum of dimensions of negative definite subspaces of smooth tests. Then
$\operatorname{index}(Q_{\varepsilon,u})
\leq\liminf_j\operatorname{index}(Q_{\varepsilon,u_j})$.
In particular a fixed upper bound for the indices of the critical points passes to the limit in Theorem \ref{thm:phase-fixed-parameter-compactness}.
\end{proposition}
\begin{proof}
Take any finite-dimensional negative subspace for the limit form and a smooth basis $\phi_1,\ldots,\phi_m$. Since $|u_j^2-u^2|\leq2|u_j-u|$, the potential terms in every matrix entry $Q_{\varepsilon,u_j}(\phi_r,\phi_s)$ converge to those of the limit form by $L^2$ convergence and boundedness of the fixed tests. The negative definite limit matrix is bounded above by $-\delta$ times the Euclidean norm squared on its coefficient vectors for some $\delta>0$, using compactness of the unit sphere in $\mathbb R^m$. Entrywise convergence preserves negative definiteness for all large $j$. Therefore the approximating indices are eventually at least $m$. Taking all possible finite $m$ proves the claim, including the case of an infinite limit index. No convergence of eigenfunction bases is assumed.
\end{proof}

\begin{proposition}[A sharp necessary parameter range for nonconstant solutions]\label{prop:phase-poincare-threshold}
Assume $M$ is connected, and let
\[
 \lambda_1=\inf_{\substack{v\in H^1(M),\ \int_Mv=0\\v\ne0}}
          \frac{\int_M|\nabla v|^2}{\int_Mv^2}>0.
\]
If $\varepsilon^2\lambda_1\geq1$, every smooth Allen--Cahn critical point is a constant in $\{-1,0,1\}$. On a circle this threshold is sharp, as the periodic examples exist for every $\varepsilon^2\lambda_1<1$.
\end{proposition}
\begin{proof}
Corollary \ref{cor:closed-manifold-poincare} gives positivity of $\lambda_1$ and the Poincare inequality.

Let $\bar u$ be the average of a critical point and set $v=u-\bar u$. Testing its equation with $v$ gives
\[
 \varepsilon^2\int|\nabla u|^2-\int v^2
    +\int(u^3-\bar u^3)(u-\bar u)=0.
\]
The first two terms together are nonnegative by $\varepsilon^2\lambda_1\geq1$ and Poincare. The final integrand is nonnegative because the cubic is strictly increasing, and vanishes only when $u=\bar u$. It follows that $u$ is constant, including in the equality case of the threshold. The equation then forces one of the three stated values.

On the circle of length $\ell$, $\lambda_1=(2\pi/\ell)^2$, from the Fourier modes. Example \ref{ex:phase-circle-periods} supplies a one-oscillation solution exactly for $\varepsilon<\ell/(2\pi)$, proving sharpness on this model. No sufficiency claim for a general manifold is made.
\end{proof}
